\section*{Capítulo \ref{ch:b1:fluid-statics} --- Estática dos fluidos}

\begin{solution}{exo:b1:fluid-statics:1}
$P = P_0 + \rho gh$: \qty{2.0e5}{Pa} (\qty{2.0}{bar}), \qty{1.1e6}{Pa}
(\qty{11}{bar}), \qty{1.0e7}{Pa} (\qty{101}{bar}). A \qty{300}{m}: $P =
\qty{3.1e6}{Pa}$, força \qty{3.1}{MN} sobre \qty{1}{m^2} (resultante, contra
\qty{1}{atm} do lado de dentro: \qty{3.0}{MN}).
\end{solution}

\begin{solution}{exo:b1:fluid-statics:2}
$h = P/\rho g$: mercúrio \qty{0.760}{m}; água \qty{10.3}{m}. Uma bomba pode,
no melhor dos casos, fazer vácuo acima da água: a atmosfera então empurra a
coluna para cima de \qty{10.3}{m}, e não mais.
\end{solution}

\begin{solution}{exo:b1:fluid-statics:3}
$F_2 = F_1S_2/S_1 = \qty{10}{kN}$ (uma tonelada); $d_2 = d_1S_1/S_2 = \qty{2.0}{mm}$
(mesmo volume deslocado, mesmo trabalho).
\end{solution}

\begin{solution}{exo:b1:fluid-statics:4}
Mesma pressão no nível da interface óleo--água nos dois ramos: $\rho_og
h_o = \rho_wgh_w$, $h_w = 0.85 \times 12 = \qty{10.2}{cm}$ de água acima
desse nível no outro ramo; a superfície do óleo fica $12 - 10.2 = \qty{1.8}{cm}$
mais alta.
\end{solution}

\begin{solution}{exo:b1:fluid-statics:5}
$H = RT/Mg = 8.314 \times 273/(0.029 \times 9.81) = \qty{8.0}{km}$; $P/P_0 =
\eu^{-z/H}$: $0.54$ a \qty{5}{km}, $0.33$ a \qty{8848}{m}; metade em $H\ln2 =
\qty{5.5}{km}$ --- os números usados no \cref{ch:b1:kinetic-theory}.
\end{solution}

\begin{solution}{exo:b1:fluid-statics:6}
$F = \tfrac12\rho gh^2L = \qty{1.96e8}{N}$; \omterm{prop:b1:fluid-statics:wall}{centro de pressão} à profundidade de \qty{13.3}{m},
isto é, \qty{6.7}{m} acima da base; momento em relação à base $F \times
6.7 = \qty{1.3e9}{N.m}$. Curvar a barragem na direção da água transforma o empurrão
em compressão do arco, levada às encostas do vale.
\end{solution}

\begin{solution}{exo:b1:fluid-statics:7}
$917/1025 = 89\%$; tronco $60\%$. Balsa: empuxo com imersão total $1000 \times
2.0 \times 9.81 = \qty{19.6}{kN}$, peso próprio \qty{11.8}{kN}: \qty{7.8}{kN}
de sobra, dez pessoas\dots{} a décima com os pés molhados: nove para ficar
seco (\qty{800}{kg}).
\end{solution}

\begin{solution}{exo:b1:fluid-statics:8}
Volume imerso $m/\rho$: em água \qty{20}{cm^3}; em salmoura $20/1.1 = \qty{18.2}{cm^3}$,
$\qty{1.8}{cm^3}$ a menos, isto é, $1.8/0.50 = \qty{3.6}{cm}$ a mais de haste: \qty{7.6}{cm}
emersos. Sensibilidade $\approx m/(\rho^2S) \approx \qty{40}{cm}$ por unidade de densidade
relativa (\qty{3.6}{cm} por $0.1$).
\end{solution}

\begin{solution}{exo:b1:fluid-statics:9}
Volume deslocado $2800/1.025 = \qty{2732}{m^3}$: \qty{268}{m^3} ($9\%$)
emergem. Equilíbrio indiferente: massa $1025 \times 3000 = \qty{3075}{t}$: admitir \qty{275}{t}.
Água doce: o empuxo cai para \qty{3000}{t}: \qty{75}{t} pesado demais --- ele afunda
a menos que \qty{75}{t} sejam bombeadas para fora.
\end{solution}

\begin{solution}{exo:b1:fluid-statics:10}
Ar deslocado: $10^{-2} \times 1.20 = \qty{12.0}{g}$; hélio dentro $12.0 \times
4/29 = \qty{1.7}{g}$; borracha \qty{3.0}{g}: sustentação líquida \qty{7.3}{g} --- sete metros
de barbante leve.
\end{solution}

\begin{solution}{exo:b1:fluid-statics:11}
Momento $\int_0^h\rho gx\cdot xL\,\dd x = \rho gLh^3/3$, força $\rho gLh^2/2$:
braço $2h/3$. De $h_1$ a $h_2$: momento $\rho gL(h_2^3 - h_1^3)/3$, força
$\rho gL(h_2^2 - h_1^2)/2$: profundidade $\frac23(h_2^3 - h_1^3)/(h_2^2 - h_1^2)$.
\end{solution}

\begin{solution}{exo:b1:fluid-statics:12}
Uma imersão suplementar $x$ acrescenta o empuxo $\rho gSx$ para cima: $m\ddot x = -\rho gSx$,
$\omega_0^2 = \rho gS/m$, $T = 2\pi\sqrt{m/\rho gS} = 2\pi\sqrt{0.020/(1000 \times
9.81 \times 5\times10^{-5})} = \qty{1.3}{s}$.
\end{solution}

\begin{solution}{pb:b1:fluid-statics:1}
\textbf{1.} $P = 1.013\times10^5 + 1025 \times 9.81 \times 300 = \qty{3.12e6}{Pa}
= \qty{31}{bar}$.

\textbf{2.} $S = \pi(0.40)^2 = \qty{0.503}{m^2}$: força resultante $(P - P_0)S =
\qty{1.52e6}{N}$ --- 155 toneladas, um caminhão carregado sobre uma escotilha.

\textbf{3.} Quilha \qty{10}{m} mais fundo: $+\rho g \times 10 = \qty{1.0e5}{Pa}$ a mais,
$3\%$: o casco é carregado quase uniformemente.

\textbf{4.} $\Delta\rho/\rho = \chi_T\Delta P = 4.6\times10^{-10} \times 3.0\times10^6
= 0.14\%$: desprezível para a pressão (alguns quilopascals em três
milhões).

\textbf{5.} Trinta bars do lado de fora, um do lado de dentro: o casco é um
vaso de pressão carregado para dentro; um cilindro (ou uma esfera) transforma
a pressão em compressão ao longo da parede, coisa que o aço resiste bem,
ao passo que uma chapa plana se dobraria.

\textbf{6.} $\Delta P\,R/e = 3.0\times10^6 \times 5.0/0.030 = \qty{5.0e8}{Pa}
= \qty{500}{MPa}$: perto do limite de escoamento --- \qty{300}{m} está no
limite para este casco, e barcos mais profundos exigem aço mais espesso ou
mais resistente (ou titânio).

\textbf{7.} Volume deslocado $m_0/\rho = 2800/1.025 = \qty{2732}{m^3}$: \qty{268}{m^3}
emergem ($9\%$).

\textbf{8.} Equilíbrio indiferente: $m = \rho V = \qty{3075}{t}$: admitir \qty{275}{t}.

\textbf{9.} O empuxo cai de $0.1\%$ de $\rho Vg$: $0.001 \times 3075 \times 9.81
\times10^3 = \qty{30}{kN}$ para baixo (três toneladas). Mais fundo $\Rightarrow$ mais
compressão $\Rightarrow$ mais pesado: instável --- um submarino tem de
corrigir o lastro continuamente (e um casco de aço se comprime menos do que
a água, o que ajuda; um casco flexível demais afundaria sem volta).

\textbf{10.} Empuxo $1000 \times 3000 \times 9.81 = \qty{29.4}{MN}$ contra o
peso $3075 \times 9.81 \times 10^3 = \qty{30.2}{MN}$: \qty{0.74}{MN} para baixo
(\qty{75}{t}); bombear \qty{75}{t} para fora.

\textbf{11.} Imediatamente \qty{20}{t} mais leve: admitir \qty{20}{t} de água
(\qty{19.5}{m^3}) nos tanques de compensação.

\textbf{12.} A água só pode ser expulsa por algo a pressão mais alta do que
a do mar do lado de fora; o ar comprimido dos cilindros a empurra para fora
--- mas só enquanto a pressão dos cilindros exceder a pressão ambiente
naquela profundidade.

\textbf{13.} $PV$ constante: \qty{275}{m^3} a \qty{31}{bar} exigem $275 \times 31
= \qty{8500}{m^3}$ de ar a \qty{1}{atm} --- guardados a \qty{200}{bar} em
\qty{43}{m^3} de cilindros.

\textbf{14.} \qty{2.0}{bar} a \qty{10}{m}; \qty{4.0}{bar} a \qty{30}{m}.

\textbf{15.} $PV$ constante: $6.0/4.0 = \qty{1.5}{L}$ --- a caixa torácica é
esmagada para dentro; os apneístas sentem isso.

\textbf{16.} O ar tomado a \qty{4}{bar} se expande quatro vezes: \qty{24}{L} em pulmões
que comportam \qty{6}{L}: ruptura. Nunca prenda a respiração na subida;
expire continuamente.

\textbf{17.} $\Delta P = \rho g \times 1.0 = \qty{1.0e4}{Pa}$; força $\qty{1.0e4}{Pa}
\times 0.05 = \qty{500}{N}$ sobre o peito: os músculos respiratórios não
levantam cinquenta quilos --- um snorkel só funciona nos primeiros decímetros.

\textbf{18.} $12 \times 200 = \qty{2400}{L}$ a \qty{1}{bar}, isto é, \qty{600}{L}
a \qty{4}{bar}: \qty{30}{min}.

\textbf{19.} $V = 10^8/1025 = \qty{97600}{m^3}$.

\textbf{20.} Em água doce ele precisa deslocar $10^5\ \unit{m^3}$: \qty{2400}{m^3}
a mais, isto é, $2400/8000 = \qty{0.30}{m}$ mais fundo.

\textbf{21.} $10^4/1.025 = \qty{9756}{m^3}$ a mais: \qty{1.2}{m} mais fundo.

\textbf{22.} Vazio, o $G$ do navio fica alto e pouco casco está imerso:
o metacentro pode cair abaixo de $G$ e o navio emborcar numa onda. Água do
mar no fundo do casco baixa $G$ e aumenta o calado.

\textbf{23.} $\rho gh = \qty{2.0e5}{Pa}$, \qty{0.2}{MN/m^2}: quinze vezes menos
do que os \qty{3}{MN/m^2} do submarino.

\textbf{24.} Não: a pressão à profundidade do submarino vale $P_0 + \rho gh$,
fixada só pela profundidade. O peso do petroleiro é sustentado pela água que
ele desloca, o que eleva o nível do mar em toda parte de uma quantidade
imensurável; nenhuma carga suplementar chega ao submarino.

\textbf{25.} $\dd P/\dd z = -\rho g$ --- donde $P = P_0 + \rho gh$ para toda
profundidade, força e volume pulmonar aqui --- e o teorema de Arquimedes, que é
essa relação integrada sobre uma superfície fechada.
\end{solution}
